\chapter[Bases de Datos NoSQL]{\emj{🌐}~Bases de Datos NoSQL}

\section[¿Cuándo NoSQL?]{\emj{🧠}~¿Cuándo NoSQL?}

``NoSQL'' (Not Only SQL) es un término paraguas para sistemas de almacenamiento que no siguen el modelo relacional estricto.
No significa que SQL sea malo — significa que existen patrones de datos y acceso para los que el modelo relacional no es la herramienta más adecuada.
La elección correcta requiere entender el \emph{modelo de datos}, los \emph{patrones de acceso} y las \emph{garantías} que el sistema necesita.

El modelo relacional exige definir el esquema antes de insertar datos, enforcea integridad referencial, y usa joins para combinar datos de múltiples tablas.
Estas características son ventajas para datos estructurados con relaciones complejas y transacciones ACID.
Pero se convierten en fricción cuando: el esquema cambia frecuentemente durante el desarrollo, cada registro tiene atributos diferentes, o se necesita escalar escrituras horizontalmente a través de múltiples nodos sin coordinación.

\begin{teoriabox}{¿Cuándo considerar NoSQL?}
\begin{itemize}
  \item \textbf{Esquema variable:} un catálogo de e-commerce donde cada producto (laptop, zapato, alimento) tiene atributos completamente distintos. Un modelo documental almacena cada producto como un documento JSON con su propia estructura.
  \item \textbf{Escalado horizontal de escritura:} sistemas que reciben millones de escrituras por segundo (IoT, telemetría, logs). Cassandra puede distribuir escrituras entre cientos de nodos sin coordinación.
  \item \textbf{Caché de baja latencia:} sesiones de usuario, resultados de queries costosas, contadores en tiempo real. Redis sirve datos desde RAM con latencia de microsegundos.
  \item \textbf{Relaciones como dato principal:} redes sociales, grafos de conocimiento. En Neo4j, traversar 5 niveles de ``amigos de amigos'' es una operación nativa; en SQL requeriría CTEs recursivas costosas.
\end{itemize}
Lo que NoSQL generalmente \emph{no} resuelve mejor que SQL: transacciones complejas multi-entidad, queries ad-hoc arbitrarias, integridad referencial automática, joins flexibles.
\end{teoriabox}

\section[Tipos y cuándo usar cada uno]{\emj{📋}~Tipos y cuándo usar cada uno}

Cada categoría de NoSQL tiene un modelo de datos distinto, optimizado para un patrón de acceso específico.
Elegir una base de datos columnar para un caso de uso que necesita documentos heterogéneos, o un grafo para un caso de uso que necesita caché, produce sistemas más complicados y lentos que un RDBMS bien diseñado.

\begin{itemize}
  \item \textbf{Documental (MongoDB, CouchDB, Firestore):} cada registro es un documento JSON/BSON autocontenido, sin esquema rígido.
    Las relaciones se manejan embebiendo documentos dentro de otros (desnormalización intencional) o con referencias manuales.
    Ideal para: catálogos de productos, CMS, perfiles de usuario con atributos variables, datos que se leen y escriben como unidad.
    Limitación: joins entre colecciones son costosos y no garantizados ACID en todas las versiones.

  \item \textbf{Clave-valor (Redis, DynamoDB, Memcached):} el modelo más simple — un mapa de clave a valor opaco.
    Acceso O(1) por clave exacta. Redis enriquece el modelo con tipos de valor (strings, listas, sets, sorted sets, hashes, streams).
    Ideal para: caché de sesiones, rankings en tiempo real, colas de trabajo, contadores atómicos, pub/sub.
    Limitación: no soporta consultas por valor, solo por clave exacta.

  \item \textbf{Columnar (Cassandra, HBase, ClickHouse):} los datos se almacenan agrupados por columna, no por fila.
    Optimizado para leer muchas filas de pocas columnas (analítica) y para escritura masiva distribuida.
    Cassandra particiona datos por una partition key y los ordena por clustering keys.
    Ideal para: series de tiempo, logs de eventos, métricas de IoT, historial de actividad.
    Limitación: el modelo de consulta es rígido — diseñas las tablas en función de las queries, no al revés.

  \item \textbf{Grafos (Neo4j, Amazon Neptune, JanusGraph):} nodos (entidades) y aristas (relaciones) con propiedades en ambos.
    Traversar relaciones es O(1) por arista, independiente del tamaño total del grafo.
    Ideal para: detección de fraude (encontrar patrones en redes de transacciones), recomendaciones (``quienes compraron X también compraron Y''), gestión de identidades y acceso.
    Limitación: no escala horizontalmente tan bien como otros NoSQL; menos herramientas y ecosistema.
\end{itemize}

\section[Teorema CAP]{\emj{⚖️}~Teorema CAP}

El teorema CAP (Eric Brewer, 2000) establece una limitación fundamental de los sistemas distribuidos: ningún sistema puede garantizar simultáneamente las tres propiedades. Durante una partición de red (falla de comunicación entre nodos), el sistema debe elegir entre consistencia y disponibilidad.

\begin{itemize}
  \item \textbf{C} (Consistency — Consistencia): todos los nodos del sistema ven exactamente los mismos datos al mismo tiempo. Una escritura en el nodo A es inmediatamente visible al leer en el nodo B.
  \item \textbf{A} (Availability — Disponibilidad): el sistema siempre responde a las requests, aunque la respuesta pueda ser con datos desactualizados. Nunca retorna error por indisponibilidad.
  \item \textbf{P} (Partition Tolerance — Tolerancia a particiones): el sistema sigue funcionando aunque haya fallas de red que dividan el cluster en grupos que no pueden comunicarse entre sí.
\end{itemize}

En la práctica, la Tolerancia a Particiones (P) es \textbf{no negociable} en cualquier sistema distribuido real: las redes fallan.
Por lo tanto, la elección real es entre C y A durante una partición:

\begin{tipbox}{CAP y bases de datos reales}
\begin{itemize}
  \item \textbf{CP} (Consistency + Partition Tolerance): MongoDB, HBase, ZooKeeper. Durante una partición, prefieren rechazar requests (perder disponibilidad) antes que servir datos inconsistentes. El sistema puede estar ``temporalmente caído'' pero los datos son siempre correctos.
  \item \textbf{AP} (Availability + Partition Tolerance): Cassandra, DynamoDB, CouchDB. Durante una partición, siguen respondiendo con los datos que tienen, aunque puedan estar desactualizados. Usan \textbf{consistencia eventual}: con el tiempo, todos los nodos convergen al mismo valor.
  \item \textbf{CA} (Consistency + Availability sin Partition Tolerance): solo posible en un único nodo. En la práctica, un RDBMS en un solo servidor. En cuanto hay replicación, P está siempre presente.
\end{itemize}
La consistencia eventual significa: ``si no hay nuevas escrituras, eventualmente todos los nodos tendrán el mismo valor''. El tiempo de convergencia depende del sistema y la carga. Es aceptable cuando los usuarios pueden tolerar ver datos levemente desactualizados (timelines de redes sociales, contadores aproximados).
\end{tipbox}

\begin{lstlisting}[language=, caption={MongoDB --- CRUD básico}]
// Insertar documento
db.productos.insertOne({
  nombre: "Laptop",
  precio: 15000,
  specs: { ram: "16GB", ssd: "512GB" },
  tags: ["electrónica", "computadoras"]
});

// Consultar con filtro anidado
db.productos.find(
  { "specs.ram": "16GB", precio: { $lt: 20000 } },
  { nombre: 1, precio: 1, _id: 0 }
);

// Agregar con $group (equivale a GROUP BY)
db.ventas.aggregate([
  { $match: { fecha: { $gte: new Date("2024-01-01") } } },
  { $group: { _id: "$region", total: { $sum: "$monto" } } },
  { $sort: { total: -1 } }
]);
\end{lstlisting}

\begin{lstlisting}[language=, caption={Redis --- caché y estructuras de datos}]
# String con TTL
SET session:user:123 '{"nombre":"Ana","rol":"admin"}' EX 3600

# Contador atómico
INCR visitas:articulo:456
EXPIRE visitas:articulo:456 86400

# Lista (cola de trabajo)
LPUSH cola:emails "job:uuid-1234"
RPOP cola:emails

# Set ordenado (ranking)
ZADD ranking 1500 "usuario:ana"
ZADD ranking 2300 "usuario:carlos"
ZREVRANGE ranking 0 9 WITHSCORES
\end{lstlisting}

\begin{entrevistabox}
\begin{enumerate}
  \item ¿Cuándo elegiría MongoDB sobre PostgreSQL para almacenar datos de productos?
  \item Explica el teorema CAP y dónde encajan Redis y Cassandra.
  \item ¿Cuál es la diferencia entre consistencia fuerte y eventual?
\end{enumerate}
\end{entrevistabox}
